\section{Problem setting}
\label{sec:setting}

\paragraph{Channels.} The specification covers six channels: \texttt{MOBILE\_MONEY},
\texttt{USSD}, \texttt{AGENT\_BANKING}, \texttt{CARD}, \texttt{ONLINE} and
\texttt{BANK\_TRANSFER}. They differ in what a fraud system can observe, and those differences are
built into the feature contract rather than handled as missing data after the fact.

\paragraph{USSD.} USSD sessions run on basic handsets and carry no device fingerprint. Every
transaction still yields a feature vector with the same slots; the four features derived from a
device fingerprint are structurally missing (NaN) for USSD and handled by the models' native
missing-value logic, so the remaining features are computed from data (resolution D-04). The
parity suite requires the NaN positions to match exactly on both computation paths
(Section~\ref{sec:features}), because a NaN on one path and a number on the other is a contract
violation rather than a numerical difference.

\paragraph{Agents.} Agent banking, where a registered agent handles cash-in and cash-out on a
customer's behalf, has its own fraud (float manipulation, fake cash-outs) and its own signals: the
agent's float utilisation, cash-out counts, distinct customers and distance from the registered
location. Those features are defined only for \texttt{AGENT\_BANKING} and are NaN on every other
channel, again by contract.

\paragraph{Corridors.} Transfers within a country, within a regional bloc such as the East African
Community (EAC), and beyond it are different populations: regional remittances, including
cross-border salary and family transfers, are routine and must not be treated like arbitrary
foreign transfers. The categorical \texttt{corridor\_class} feature has four classes, derived from
the two countries' bloc memberships rather than from a list of EAC countries (ADR 0023):
\texttt{DOMESTIC}, \texttt{INTRA\_BLOC}, \texttt{CROSS\_BLOC\_AFRICA} and
\texttt{INTERCONTINENTAL}. \textbf{Only the first two occur on this benchmark}: every simulated
counterparty is in one of the five countries, all of them EAC members, so the models never saw a
cross-bloc or intercontinental transfer and nothing here says how they would score one (the
feature registry records the two classes as unreachable, PB-43).

\paragraph{Currencies, time and money.} Five currencies with different minor units, amounts
normalised to Rwandan francs by a dated exchange-rate table, local time derived from coordinates
(the Democratic Republic of the Congo spans two time zones), and amounts carried as decimals,
never binary floating point.

\paragraph{Decision path.} The system returns a decision (approve, decline, or hold for review)
synchronously to core banking, with a server-side latency target of 40~ms at the 95th percentile
for flagged transactions (the specification's latency budget). A decision that holds a transaction for an
analyst is resolved by webhook, not by an open request. Scores are calibrated probabilities from
a weighted XGBoost and LightGBM ensemble, and every flagged decision carries an exact TreeSHAP
explanation. The scoring service and decision engine are milestones M5 and M6; their measured
latency is \notyet{M5, M6 and M10}.

\paragraph{What is not in scope.} Graph-based features (the specification's later phases),
real customer data of any kind, and any claim about a real institution, regulator or deployment.
